\documentclass[11pt]{article}
\usepackage[margin=1in]{geometry}
\usepackage{amsmath,amssymb,amsthm}
\usepackage{hyperref}
\usepackage{enumitem}
\setlist{nosep,leftmargin=*}
\newtheorem*{question}{Question}
\title{Working Notes: Fruitfulness, Compression,\\and the Shape of Mathematics}
\author{Vasily Shablov \and with Claude (Anthropic)}
\date{October 2026 --- current thoughts, not a paper}
\begin{document}
\maketitle

\begin{quote}\small
``Civilization advances by extending the number of important operations which we can
perform without thinking about them.'' \hfill --- A.\,N.~Whitehead, 1911
\end{quote}

\section{Where we are}

We asked an open question, \emph{what is the fruitfulness of a mathematical idea?}, and gave
it a precise answer.

\paragraph{The definition.} A named concept (a definition or a lemma) is a
\emph{compression}: its body is written once and cited afterwards.

\paragraph{What is proved in Lean.} No \texttt{sorry} anywhere.
\begin{itemize}
\item Naming a body of size $b$ that is used $k$ times saves exactly $(k-1)(b-1)-1$ symbols.
\item Unfolding every citation at most doubles the size per result.
\item That doubling is attained, and a Fibonacci library grows by the golden ratio per result.
\end{itemize}

\paragraph{What we measured.} We exported all 771{,}129 constants of Mathlib, with
20.4\,M citation edges, and analysed the 315{,}085 user-facing ones.
\begin{itemize}
\item \textbf{Naive unfolding explodes.} Unfolded size grows about $1.5\times$ per
  dependency level ($R^2=0.98$). At the deepest level it reaches roughly $10^{65.6}$ symbols.
\item \textbf{Reuse rediscovers the textbook.} Ranked by reuse, the top concepts are
  \texttt{Set}, \texttt{Category}, \texttt{Functor}, \texttt{TopologicalSpace},
  \texttt{Module}, and so on.
\item \textbf{Two axes of human taste, each computable.}
  \begin{itemize}
  \item Concepts that humans call important are predicted by \emph{reuse}: AUC $0.87$, or
    $0.86$ when compared only within the same file.
  \item Theorems that humans call famous are predicted by \emph{depth}: AUC $0.78$, or $0.66$
    within the same file.
  \end{itemize}
\end{itemize}

\noindent These notes record what I now think these results \emph{mean}. Most of it is
questions.

\section{A correction I owe first}

We said ``fully unfolded, Mathlib does not fit in the universe.'' That is true of
\emph{naive} inlining. But $U$ is only an \emph{upper} bound on the size of a lemma-free
proof. A clever re-prover can restructure a proof instead of copying lemma bodies.

For Frege systems this is a theorem: tree-like Frege polynomially simulates dag-like Frege
(Kraj\'i\v{c}ek). There, reusing lemmas buys at most a \emph{polynomial} saving, even though
naive unfolding is exponential. Our exponential theorems are correct, but they are about
inlining, not about the minimal lemma-free proof. The honest claim is:
\begin{itemize}
\item \emph{lemma names} make mathematics exponentially shorter than \emph{copying};
\item whether \emph{some} lemma-free proof would be short is a question of
  \textbf{proof complexity}, and it is largely open.
\end{itemize}

\section{The deepest form of our question is a famous open problem}

There are two kinds of naming, and they are not the same.

\begin{enumerate}
\item \textbf{Lemmas} name a \emph{proved statement}: reuse a result. In proof complexity
  this is dag-like versus tree-like proofs.
  \begin{itemize}
  \item For resolution, the gap is exponential (Ben-Sasson, Impagliazzo and Wigderson).
  \item For Frege, the gap is polynomial (Kraj\'i\v{c}ek).
  \end{itemize}
\item \textbf{Definitions} introduce a \emph{new symbol} that abbreviates a formula. That
  is Tseitin's extension rule, and Extended Frege is Frege plus definitions.
  \begin{itemize}
  \item Whether Extended Frege is super-polynomially stronger than Frege is one of the
    central open problems of proof complexity.
  \end{itemize}
\end{enumerate}

\begin{question}[Q1, sharpened]
Are \emph{new concepts} ever exponentially necessary, beyond reusing lemmas? Or can every
proof that uses definitions be rewritten, at polynomial cost, using only old vocabulary?
\end{question}

\noindent This is exactly the question ``are human concepts forced or merely convenient?'',
in formal dress.

Going up a type level, the answer is spectacular. G\"odel (1936) showed that moving to
higher-order logic shortens some proofs by more than any computable function. Boolos's
``curious inference'' (1987) is a concrete case: a valid first-order argument whose shortest
first-order proof is unimaginably long, physically impossible to write down, while a short
second-order proof exists. So at the level of \emph{types of concepts}, concepts are provably
forced. At the propositional level, nobody knows.

\begin{question}[Q1$'$]
Mathlib is a vast empirical sample of definitions in use. Can its data say anything about
\emph{where} definitions give super-polynomial savings? That would give experimental
evidence on a problem that has resisted proof.
\end{question}

\section{Value has two axes: depth and leverage}

Our human-label results split cleanly.
\begin{itemize}
\item \textbf{Famous theorems are deep}, in the sense of high $\log U - \log s$. They
  certify a large amount of hidden work behind a short statement.
\item \textbf{Fundamental concepts have leverage.} They are small and reused enormously.
\end{itemize}

Both axes already exist in the theory of information, independently of mathematics.
\begin{itemize}
\item \textbf{Bennett's logical depth} (1988). An object is valuable when it is the result
  of a long computation that cannot be shortcut: a genome, a solved theorem. Bennett
  proposed depth, not information content, as the measure of value. Our finding that fame
  tracks unfolded depth is, as far as I know, the first large-scale test of Bennett's
  thesis on human judgments of mathematical value.
\item \textbf{Schmidhuber's compression progress} (2009). An observer finds something
  interesting when it improves the observer's compressor; the first derivative of
  compressibility. Our $(k-1)(b-1)-1$ is exactly the compression progress a single new
  concept produces. Reuse predicting human concept taste (0.87) is evidence for
  Schmidhuber's thesis in mathematics.
\end{itemize}

\begin{question}[Q2]
Are depth and leverage the \emph{only} two axes of mathematical value? What explains the
residual, the roughly 15--20\% of human judgments neither measure predicts? Is it surprise,
beauty, connection between distant fields, or just historical accident? Can a third
measure, such as \emph{bridging} (citations spanning distant parts of the library), absorb
it?
\end{question}

\noindent This bears directly on the debate that started this. ``Nothing to optimize'' is
false in a strong sense. Theories of the \emph{reward} for open-ended inquiry exist
(Bennett, Schmidhuber), and here they predict human taste.

\section{A glue law across nature}

To see meaning in Mathlib we had to remove a layer of \emph{typeclass plumbing}:
\begin{itemize}
\item instance lemmas that are used thousands of times;
\item each one carries almost no content;
\item no human would ever call one important.
\end{itemize}
The same layer appears everywhere reuse appears.
\begin{itemize}
\item \textbf{Language.} The most frequent words are \emph{function words} (\emph{the},
  \emph{of}). They are pure glue, and corpus linguists filter them out as ``stop words''.
\item \textbf{Metabolism.} The hubs of metabolic networks are \emph{currency metabolites}
  (ATP, NADH, water). Network biologists must remove them before pathways become visible
  (Jeong et al.~2000; Ma and Zeng~2003).
\item \textbf{Software.} The most depended-upon packages are tiny utilities. The
  \texttt{left-pad} incident is the famous case.
\item \textbf{Mathlib.} Coercions, instance projections, \texttt{mul\_one}.
\end{itemize}

There are more parallels.
\begin{itemize}
\item \textbf{Heavy-tailed reuse} appears in all of these. Mathlib has
  $P(k\ge x)\propto x^{-0.8}$, with 165 constants used more than 10{,}000 times. Compare
  Zipf's law for words, protein-domain family sizes (Koonin et al.), and citation counts
  (Price's cumulative advantage).
\item \textbf{The law of abbreviation} holds as well. In natural language, frequent words
  are shorter (Zipf). We just checked Mathlib, and it holds there too.
  \begin{itemize}
  \item The mean name length of definitions falls monotonically with use: from 13.4
    characters (never used) to 7.1 (used $10^4$--$10^5$ times) to 3 (\texttt{Set},
    \texttt{coe}).
  \item Spearman $\rho(\text{uses},\text{name length})=-0.21$.
  \item Spearman $\rho(\text{uses},\text{body size})=-0.25$.
  \end{itemize}
  Nobody designed this. Mathlib's naming convention derives names from statements, so the
  law emerges from which concepts get reused.
\end{itemize}

\begin{question}[Q3: the glue theorem]
Is there a single theorem behind all of these? Consider any system that grows by reuse
under a cost on description length. Must it develop (i) heavy-tailed reuse, (ii) a
high-frequency, low-content glue layer, and (iii) abbreviation of the frequent?
Preferential attachment explains (i), and Zipf's least-effort argument gestures at (iii).
A unified compression-theoretic derivation of all three would connect linguistics,
metabolism, software and mathematics.
\end{question}

\paragraph{Why Mathlib matters here.} It is the \emph{model organism} of reuse systems.
\begin{itemize}
\item In biology, edges are inferred and noisy. In language, ``concept'' is fuzzy.
\item In Mathlib, every edge is \emph{exact} and every unit is \emph{checked}.
\end{itemize}
That makes laws testable here first, with unusual precision.

\section{Hierarchy, decoupling, and why the world is compressible}

\paragraph{Simon's watchmakers.} Herbert Simon's \emph{Architecture of Complexity} (1962)
tells of two watchmakers.
\begin{itemize}
\item One assembles stable sub-assemblies. The other assembles parts one by one.
\item Interrupted often, only the first finishes. The hierarchical builder has an
  exponential advantage.
\item Simon argued this is why nature is hierarchical: cells, organs, organisms,
  organizations.
\end{itemize}
Our Lean theorems are the \emph{size} counterpart of his \emph{time} argument.

\paragraph{Proof irrelevance is scale separation.} In Lean, any two proofs of the same
proposition are equal. A later result can depend only on the \emph{statement} of a lemma,
never on its proof. That is precisely the physicist's \emph{decoupling of scales}.
\begin{itemize}
\item Chemistry needs the periodic table, not quantum chromodynamics.
\item Renormalization explains why the microscopic details wash out.
\item The quantity $\log U-\log s$ measures how much ``microstate'' a statement
  ``macrostate'' hides: an entropy of forgetting.
\end{itemize}

\begin{question}[Q4: Wigner's puzzle, as compression]
Why is mathematics unreasonably effective in physics? A conjecture: because both are
products of hierarchical compression under scale separation, so the concepts that compress
mathematics also compress the world. This is testable. Formalized physics libraries now
exist (\emph{PhysLean}). Does physics reuse the same top concepts as Mathlib, and with the
same heavy tail and glue layer?
\end{question}

\section{The growth constant}

Naive unfolding grows by a factor of about $1.5$ per level:
\begin{itemize}
\item $1.45$, $1.64$ and $1.41$ across the three thirds of the height range;
\item for comparison, the golden ratio $1.618$ for a library citing two predecessors, and
  $2$ for the provable worst case.
\end{itemize}
Read as a branching process, mathematics is \emph{supercritical}: each level hides about
1.5 times the work of the previous one.

\begin{question}[Q5]
\begin{itemize}
\item Is this constant a property of mathematics, or of Mathlib's style?
\item Is it stable over Mathlib's history, and the same in Isabelle or Rocq?
\item Does it relate to the ``effective number of essential citations'', via the
  $d$-bonacci constants (golden ratio for $d=2$, approaching $2$)?
\end{itemize}
\end{question}

\section{Delayed reward, measured}

The strongest remaining point of the opposing argument: fruitfulness is visible only
\emph{decades} later (Lobachevsky), so it cannot be trained on. This is now an experiment.
\begin{itemize}
\item Mathlib's git history records when every declaration was born and when each use
  appeared.
\item If a concept's reuse after a few months predicts its reuse years later (cumulative
  advantage), then the effective reward horizon is months, not decades.
\item If it does not, the delayed-reward objection gains real empirical support.
\end{itemize}
Either outcome is informative. I would like to know the answer more than I would like to
win the argument.

\section{Understanding is conditional compression}

Thurston said mathematics is about human understanding. Here is a way to make that precise.
\begin{itemize}
\item An agent $h$ \emph{understands} a proof $\pi$ when $\pi$'s size, \emph{relative to
  $h$'s library of named concepts}, fits within $h$'s working capacity.
\item This is conditional Kolmogorov complexity, with a library in place of a universal
  machine.
\item It also matches the cognitive science of chunking: chess masters see positions in
  chunks (Chase and Simon), and memory holds about seven chunks (Miller).
\end{itemize}

This definition can be plugged into the \texttt{Understands} predicate of our earlier paper.
It then makes the ``fuzzy boundary'' concrete. A proof incomprehensible to a human may be
comprehensible \emph{given a richer library}. So the question ``can humans understand AI
mathematics?'' becomes ``can the AI's library be taught?'', which is quantitative.

\section{Plan (fast, in order of value per hour)}

\begin{enumerate}
\item \textbf{Lean: the sharp $d$-bonacci bound.} This is conjecture~4 of the README.
  \begin{itemize}
  \item With at most $d$ citations, $U_i\le S\,T_i$, where $T_i=1+T_{i-1}+\dots+T_{i-d}$.
  \item It is attained by the library citing its $d$ predecessors.
  \item Proof idea: the $d$ largest earlier values of a monotone sequence are the last $d$
    values. A short exchange argument on finite sets proves it.
  \end{itemize}
\item \textbf{Formal paper} with all of the above: theorems, measurements, the correction,
  and the connections, stated as conjectures where unproved.
\item \textbf{Glue and abbreviation analysis}, made rigorous. Classify plumbing by
  \emph{kind} (instances), not by name, and report the law of abbreviation properly.
\item \textbf{Later:} the git-history experiment (Q on delayed reward), and a Stitch-style
  library-learning rediscovery test on a small part of Mathlib (Q1, ``forced vs chosen'').
\end{enumerate}

\begin{thebibliography}{99}\small
\bibitem{bss} E.~Ben-Sasson, R.~Impagliazzo, A.~Wigderson. Near optimal separation of tree-like and general resolution. \emph{Combinatorica} 24 (2004).
\bibitem{bennett} C.~H.~Bennett. Logical depth and physical complexity. In \emph{The Universal Turing Machine: A Half-Century Survey}, 1988.
\bibitem{boolos} G.~Boolos. A curious inference. \emph{J.\ Philosophical Logic} 16 (1987).
\bibitem{chase} W.~G.~Chase, H.~A.~Simon. Perception in chess. \emph{Cognitive Psychology} 4 (1973).
\bibitem{godel} K.~G\"odel. \"Uber die L\"ange von Beweisen. \emph{Ergebnisse eines math.\ Kolloquiums} 7 (1936).
\bibitem{jeong} H.~Jeong et al. The large-scale organization of metabolic networks. \emph{Nature} 407 (2000).
\bibitem{koonin} E.~V.~Koonin, Y.~I.~Wolf, G.~P.~Karev. The structure of the protein universe and genome evolution. \emph{Nature} 420 (2002).
\bibitem{krajicek} J.~Kraj\'i\v{c}ek. \emph{Bounded Arithmetic, Propositional Logic, and Complexity Theory}. CUP, 1995.
\bibitem{ma} H.~Ma, A.-P.~Zeng. Reconstruction of metabolic networks from genome data. \emph{Bioinformatics} 19 (2003).
\bibitem{miller} G.~A.~Miller. The magical number seven, plus or minus two. \emph{Psychological Review} 63 (1956).
\bibitem{price} D.~de~Solla Price. A general theory of bibliometric and other cumulative advantage processes. \emph{JASIS} 27 (1976).
\bibitem{schmidhuber} J.~Schmidhuber. Driven by compression progress. In \emph{Anticipatory Behavior in Adaptive Learning Systems}, 2009.
\bibitem{simon} H.~A.~Simon. The architecture of complexity. \emph{Proc.\ Amer.\ Phil.\ Soc.} 106 (1962).
\bibitem{thurston} W.~P.~Thurston. On proof and progress in mathematics. \emph{Bull.\ AMS} 30 (1994).
\bibitem{wigner} E.~Wigner. The unreasonable effectiveness of mathematics in the natural sciences. \emph{CPAM} 13 (1960).
\bibitem{zipf} G.~K.~Zipf. \emph{Human Behavior and the Principle of Least Effort}. 1949.
\end{thebibliography}
\end{document}Our finding that fame tracks unfolded depth reads as a large-scale test of Bennett's
  thesis on human judgments of mathematical value. I know of no prior test like it, but that
  still needs a literature check.t]{article}
\usepackage[margin=1in]{geometry}
\usepackage{amsmath,amssymb,amsthm}
\usepackage{hyperref}
\usepackage{enumitem}
\setlist{nosep,leftmargin=*}
\newtheorem*{question}{Question}
\title{Working Notes: Fruitfulness, Compression,\\and the Shape of Mathematics}
\author{Vasily Shablov \and with Claude (Anthropic)}
\date{October 2026 --- current thoughts, not a paper}
\begin{document}
\maketitle

\begin{quote}\small
``Civilization advances by extending the number of important operations which we can
perform without thinking about them.'' \hfill --- A.\,N.~Whitehead, 1911
\end{quote}

\section{Where we are}

We asked an open question, \emph{what is the fruitfulness of a mathematical idea?}, and gave
it a precise answer.

\paragraph{The definition.} A named concept (a definition or a lemma) is a
\emph{compression}: its body is written once and cited afterwards.

\paragraph{What is proved in Lean.} No \texttt{sorry} anywhere.
\begin{itemize}
\item Naming a body of size $b$ that is used $k$ times saves exactly $(k-1)(b-1)-1$ symbols.
\item Unfolding every citation at most doubles the size per result.
\item That doubling is attained, and a Fibonacci library grows by the golden ratio per result.
\end{itemize}

\paragraph{What we measured.} We exported all 771{,}129 constants of Mathlib, with
20.4\,M citation edges, and analysed the 315{,}085 user-facing ones.
\begin{itemize}
\item \textbf{Naive unfolding explodes.} Unfolded size grows about $1.5\times$ per
  dependency level ($R^2=0.98$). At the deepest level it reaches roughly $10^{65.6}$ symbols.
\item \textbf{Reuse rediscovers the textbook.} Ranked by reuse, the top concepts are
  \texttt{Set}, \texttt{Category}, \texttt{Functor}, \texttt{TopologicalSpace},
  \texttt{Module}, and so on.
\item \textbf{Two axes of human taste, each computable.}
  \begin{itemize}
  \item Concepts that humans call important are predicted by \emph{reuse}: AUC $0.87$, or
    $0.86$ when compared only within the same file.
  \item Theorems that humans call famous are predicted by \emph{depth}: AUC $0.78$, or $0.66$
    within the same file.
  \end{itemize}
\end{itemize}

\noindent These notes record what I now think these results \emph{mean}. Most of it is
questions.

\section{A correction I owe first}

We said ``fully unfolded, Mathlib does not fit in the universe.'' That is true of
\emph{naive} inlining. But $U$ is only an \emph{upper} bound on the size of a lemma-free
proof. A clever re-prover can restructure a proof instead of copying lemma bodies.

For Frege systems this is a theorem: tree-like Frege polynomially simulates dag-like Frege
(Kraj\'i\v{c}ek). There, reusing lemmas buys at most a \emph{polynomial} saving, even though
naive unfolding is exponential. Our exponential theorems are correct, but they are about
inlining, not about the minimal lemma-free proof. The honest claim is:
\begin{itemize}
\item \emph{lemma names} make mathematics exponentially shorter than \emph{copying};
\item whether \emph{some} lemma-free proof would be short is a question of
  \textbf{proof complexity}, and it is largely open.
\end{itemize}

\section{The deepest form of our question is a famous open problem}

There are two kinds of naming, and they are not the same.

\begin{enumerate}
\item \textbf{Lemmas} name a \emph{proved statement}: reuse a result. In proof complexity
  this is dag-like versus tree-like proofs.
  \begin{itemize}
  \item For resolution, the gap is exponential (Ben-Sasson, Impagliazzo and Wigderson).
  \item For Frege, the gap is polynomial (Kraj\'i\v{c}ek).
  \end{itemize}
\item \textbf{Definitions} introduce a \emph{new symbol} that abbreviates a formula. That
  is Tseitin's extension rule, and Extended Frege is Frege plus definitions.
  \begin{itemize}
  \item Whether Extended Frege is super-polynomially stronger than Frege is one of the
    central open problems of proof complexity.
  \end{itemize}
\end{enumerate}

\begin{question}[Q1, sharpened]
Are \emph{new concepts} ever exponentially necessary, beyond reusing lemmas? Or can every
proof that uses definitions be rewritten, at polynomial cost, using only old vocabulary?
\end{question}

\noindent This is exactly the question ``are human concepts forced or merely convenient?'',
in formal dress.

Going up a type level, the answer is spectacular. G\"odel (1936) showed that moving to
higher-order logic shortens some proofs by more than any computable function. Boolos's
``curious inference'' (1987) is a concrete case: a valid first-order argument whose shortest
first-order proof is unimaginably long, physically impossible to write down, while a short
second-order proof exists. So at the level of \emph{types of concepts}, concepts are provably
forced. At the propositional level, nobody knows.

\begin{question}[Q1$'$]
Mathlib is a vast empirical sample of definitions in use. Can its data say anything about
\emph{where} definitions give super-polynomial savings? That would give experimental
evidence on a problem that has resisted proof.
\end{question}

\section{Value has two axes: depth and leverage}

Our human-label results split cleanly.
\begin{itemize}
\item \textbf{Famous theorems are deep}, in the sense of high $\log U - \log s$. They
  certify a large amount of hidden work behind a short statement.
\item \textbf{Fundamental concepts have leverage.} They are small and reused enormously.
\end{itemize}

Both axes already exist in the theory of information, independently of mathematics.
\begin{itemize}
\item \textbf{Bennett's logical depth} (1988). An object is valuable when it is the result
  of a long computation that cannot be shortcut: a genome, a solved theorem. Bennett
  proposed depth, not information content, as the measure of value. Our finding that fame
  tracks unfolded depth is, as far as I know, the first large-scale test of Bennett's
  thesis on human judgments of mathematical value.
\item \textbf{Schmidhuber's compression progress} (2009). An observer finds something
  interesting when it improves the observer's compressor; the first derivative of
  compressibility. Our $(k-1)(b-1)-1$ is exactly the compression progress a single new
  concept produces. Reuse predicting human concept taste (0.87) is evidence for
  Schmidhuber's thesis in mathematics.
\end{itemize}

\begin{question}[Q2]
Are depth and leverage the \emph{only} two axes of mathematical value? What explains the
residual, the roughly 15--20\% of human judgments neither measure predicts? Is it surprise,
beauty, connection between distant fields, or just historical accident? Can a third
measure, such as \emph{bridging} (citations spanning distant parts of the library), absorb
it?
\end{question}

\noindent This bears directly on the debate that started this. ``Nothing to optimize'' is
false in a strong sense. Theories of the \emph{reward} for open-ended inquiry exist
(Bennett, Schmidhuber), and here they predict human taste.

\section{A glue law across nature}

To see meaning in Mathlib we had to remove a layer of \emph{typeclass plumbing}:
\begin{itemize}
\item instance lemmas that are used thousands of times;
\item each one carries almost no content;
\item no human would ever call one important.
\end{itemize}
The same layer appears everywhere reuse appears.
\begin{itemize}
\item \textbf{Language.} The most frequent words are \emph{function words} (\emph{the},
  \emph{of}). They are pure glue, and corpus linguists filter them out as ``stop words''.
\item \textbf{Metabolism.} The hubs of metabolic networks are \emph{currency metabolites}
  (ATP, NADH, water). Network biologists must remove them before pathways become visible
  (Jeong et al.~2000; Ma and Zeng~2003).
\item \textbf{Software.} The most depended-upon packages are tiny utilities. The
  \texttt{left-pad} incident is the famous case.
\item \textbf{Mathlib.} Coercions, instance projections, \texttt{mul\_one}.
\end{itemize}

There are more parallels.
\begin{itemize}
\item \textbf{Heavy-tailed reuse} appears in all of these. Mathlib has
  $P(k\ge x)\propto x^{-0.8}$, with 165 constants used more than 10{,}000 times. Compare
  Zipf's law for words, protein-domain family sizes (Koonin et al.), and citation counts
  (Price's cumulative advantage).
\item \textbf{The law of abbreviation} holds as well. In natural language, frequent words
  are shorter (Zipf). We just checked Mathlib, and it holds there too.
  \begin{itemize}
  \item The mean name length of definitions falls monotonically with use: from 13.4
    characters (never used) to 7.1 (used $10^4$--$10^5$ times) to 3 (\texttt{Set},
    \texttt{coe}).
  \item Spearman $\rho(\text{uses},\text{name length})=-0.21$.
  \item Spearman $\rho(\text{uses},\text{body size})=-0.25$.
  \end{itemize}
  Nobody designed this. Mathlib's naming convention derives names from statements, so the
  law emerges from which concepts get reused.
\end{itemize}

\begin{question}[Q3: the glue theorem]
Is there a single theorem behind all of these? Consider any system that grows by reuse
under a cost on description length. Must it develop (i) heavy-tailed reuse, (ii) a
high-frequency, low-content glue layer, and (iii) abbreviation of the frequent?
Preferential attachment explains (i), and Zipf's least-effort argument gestures at (iii).
A unified compression-theoretic derivation of all three would connect linguistics,
metabolism, software and mathematics.
\end{question}

\paragraph{Why Mathlib matters here.} It is the \emph{model organism} of reuse systems.
\begin{itemize}
\item In biology, edges are inferred and noisy. In language, ``concept'' is fuzzy.
\item In Mathlib, every edge is \emph{exact} and every unit is \emph{checked}.
\end{itemize}
That makes laws testable here first, with unusual precision.

\section{Hierarchy, decoupling, and why the world is compressible}

\paragraph{Simon's watchmakers.} Herbert Simon's \emph{Architecture of Complexity} (1962)
tells of two watchmakers.
\begin{itemize}
\item One assembles stable sub-assemblies. The other assembles parts one by one.
\item Interrupted often, only the first finishes. The hierarchical builder has an
  exponential advantage.
\item Simon argued this is why nature is hierarchical: cells, organs, organisms,
  organizations.
\end{itemize}
Our Lean theorems are the \emph{size} counterpart of his \emph{time} argument.

\paragraph{Proof irrelevance is scale separation.} In Lean, any two proofs of the same
proposition are equal. A later result can depend only on the \emph{statement} of a lemma,
never on its proof. That is precisely the physicist's \emph{decoupling of scales}.
\begin{itemize}
\item Chemistry needs the periodic table, not quantum chromodynamics.
\item Renormalization explains why the microscopic details wash out.
\item The quantity $\log U-\log s$ measures how much ``microstate'' a statement
  ``macrostate'' hides: an entropy of forgetting.
\end{itemize}

\begin{question}[Q4: Wigner's puzzle, as compression]
Why is mathematics unreasonably effective in physics? A conjecture: because both are
products of hierarchical compression under scale separation, so the concepts that compress
mathematics also compress the world. This is testable. Formalized physics libraries now
exist (\emph{PhysLean}). Does physics reuse the same top concepts as Mathlib, and with the
same heavy tail and glue layer?
\end{question}

\section{The growth constant}

Naive unfolding grows by a factor of about $1.5$ per level:
\begin{itemize}
\item $1.45$, $1.64$ and $1.41$ across the three thirds of the height range;
\item for comparison, the golden ratio $1.618$ for a library citing two predecessors, and
  $2$ for the provable worst case.
\end{itemize}
Read as a branching process, mathematics is \emph{supercritical}: each level hides about
1.5 times the work of the previous one.

\begin{question}[Q5]
\begin{itemize}
\item Is this constant a property of mathematics, or of Mathlib's style?
\item Is it stable over Mathlib's history, and the same in Isabelle or Rocq?
\item Does it relate to the ``effective number of essential citations'', via the
  $d$-bonacci constants (golden ratio for $d=2$, approaching $2$)?
\end{itemize}
\end{question}

\section{Delayed reward, measured}

The strongest remaining point of the opposing argument: fruitfulness is visible only
\emph{decades} later (Lobachevsky), so it cannot be trained on. This is now an experiment.
\begin{itemize}
\item Mathlib's git history records when every declaration was born and when each use
  appeared.
\item If a concept's reuse after a few months predicts its reuse years later (cumulative
  advantage), then the effective reward horizon is months, not decades.
\item If it does not, the delayed-reward objection gains real empirical support.
\end{itemize}
Either outcome is informative. I would like to know the answer more than I would like to
win the argument.

\section{Understanding is conditional compression}

Thurston said mathematics is about human understanding. Here is a way to make that precise.
\begin{itemize}
\item An agent $h$ \emph{understands} a proof $\pi$ when $\pi$'s size, \emph{relative to
  $h$'s library of named concepts}, fits within $h$'s working capacity.
\item This is conditional Kolmogorov complexity, with a library in place of a universal
  machine.
\item It also matches the cognitive science of chunking: chess masters see positions in
  chunks (Chase and Simon), and memory holds about seven chunks (Miller).
\end{itemize}

This definition can be plugged into the \texttt{Understands} predicate of our earlier paper.
It then makes the ``fuzzy boundary'' concrete. A proof incomprehensible to a human may be
comprehensible \emph{given a richer library}. So the question ``can humans understand AI
mathematics?'' becomes ``can the AI's library be taught?'', which is quantitative.

\section{Plan (fast, in order of value per hour)}

\begin{enumerate}
\item \textbf{Lean: the sharp $d$-bonacci bound.} This is conjecture~4 of the README.
  \begin{itemize}
  \item With at most $d$ citations, $U_i\le S\,T_i$, where $T_i=1+T_{i-1}+\dots+T_{i-d}$.
  \item It is attained by the library citing its $d$ predecessors.
  \item Proof idea: the $d$ largest earlier values of a monotone sequence are the last $d$
    values. A short exchange argument on finite sets proves it.
  \end{itemize}
\item \textbf{Formal paper} with all of the above: theorems, measurements, the correction,
  and the connections, stated as conjectures where unproved.
\item \textbf{Glue and abbreviation analysis}, made rigorous. Classify plumbing by
  \emph{kind} (instances), not by name, and report the law of abbreviation properly.
\item \textbf{Later:} the git-history experiment (Q on delayed reward), and a Stitch-style
  library-learning rediscovery test on a small part of Mathlib (Q1, ``forced vs chosen'').
\end{enumerate}

\begin{thebibliography}{99}\small
\bibitem{bss} E.~Ben-Sasson, R.~Impagliazzo, A.~Wigderson. Near optimal separation of tree-like and general resolution. \emph{Combinatorica} 24 (2004).
\bibitem{bennett} C.~H.~Bennett. Logical depth and physical complexity. In \emph{The Universal Turing Machine: A Half-Century Survey}, 1988.
\bibitem{boolos} G.~Boolos. A curious inference. \emph{J.\ Philosophical Logic} 16 (1987).
\bibitem{chase} W.~G.~Chase, H.~A.~Simon. Perception in chess. \emph{Cognitive Psychology} 4 (1973).
\bibitem{godel} K.~G\"odel. \"Uber die L\"ange von Beweisen. \emph{Ergebnisse eines math.\ Kolloquiums} 7 (1936).
\bibitem{jeong} H.~Jeong et al. The large-scale organization of metabolic networks. \emph{Nature} 407 (2000).
\bibitem{koonin} E.~V.~Koonin, Y.~I.~Wolf, G.~P.~Karev. The structure of the protein universe and genome evolution. \emph{Nature} 420 (2002).
\bibitem{krajicek} J.~Kraj\'i\v{c}ek. \emph{Bounded Arithmetic, Propositional Logic, and Complexity Theory}. CUP, 1995.
\bibitem{ma} H.~Ma, A.-P.~Zeng. Reconstruction of metabolic networks from genome data. \emph{Bioinformatics} 19 (2003).
\bibitem{miller} G.~A.~Miller. The magical number seven, plus or minus two. \emph{Psychological Review} 63 (1956).
\bibitem{price} D.~de~Solla Price. A general theory of bibliometric and other cumulative advantage processes. \emph{JASIS} 27 (1976).
\bibitem{schmidhuber} J.~Schmidhuber. Driven by compression progress. In \emph{Anticipatory Behavior in Adaptive Learning Systems}, 2009.
\bibitem{simon} H.~A.~Simon. The architecture of complexity. \emph{Proc.\ Amer.\ Phil.\ Soc.} 106 (1962).
\bibitem{thurston} W.~P.~Thurston. On proof and progress in mathematics. \emph{Bull.\ AMS} 30 (1994).
\bibitem{wigner} E.~Wigner. The unreasonable effectiveness of mathematics in the natural sciences. \emph{CPAM} 13 (1960).
\bibitem{zipf} G.~K.~Zipf. \emph{Human Behavior and the Principle of Least Effort}. 1949.
\end{thebibliography}
\end{document}
